\documentclass[11pt]{article}
\usepackage{amsmath,amssymb,amsthm}
\usepackage[margin=1in]{geometry}

\newtheorem{theorem}{Theorem}
\newtheorem{lemma}{Lemma}
\newtheorem{definition}{Definition}
\newtheorem{remark}{Remark}

\begin{document}

\title{Zero Topological Entropy from Regularity of Distributional Solutions}
\author{}
\date{}
\maketitle

\section{Definitions and conventions}

Let \(M\) be a compact connected smooth manifold, and let \(X\in \mathfrak X^\infty(M)\). Let \(g_t\) be the smooth flow generated by \(X\). We use the standard identification of a finite Borel measure \(\mu\) with the distribution
\[
C_c^\infty(M)\ni \varphi \longmapsto \int_M \varphi\,d\mu .
\]
Likewise, after fixing a positive smooth volume density \(\nu\) on \(M\), a smooth function \(h\in C^\infty(M)\) is identified with the distribution
\[
C_c^\infty(M)\ni \varphi \longmapsto \int_M h\varphi\,d\nu .
\]
Changing \(\nu\) by a smooth positive factor does not affect whether a distribution is represented by a smooth function.

For a distribution \(u\), the Lie derivative \(L_Xu\) is defined by
\[
\langle L_Xu,\varphi\rangle
=
-\langle u, X\varphi\rangle,
\qquad \varphi\in C_c^\infty(M).
\]
Thus, for a finite Borel measure \(\mu\),
\[
\langle L_X\mu,\varphi\rangle
=
-\int_M X\varphi\,d\mu
=
-\left.\frac{d}{dt}\right|_{t=0}\int_M \varphi(g_t x)\,d\mu(x).
\]
Hence
\[
L_X\mu=0
\]
if and only if \(\mu\) is invariant under the flow \(g_t\).

The hypothesis on \(X\) is:

\begin{quote}
If \(f\in C^\infty(M)\), \(u\) is a distribution, and \(L_Xu=f\), then \(u\) is represented by a smooth function.
\end{quote}

\section{Periodic orbits give nonsmooth invariant distributions}

\begin{lemma}
Let \(\gamma\) be a periodic orbit of \(X\). If \(\gamma\neq M\), then the normalized orbital measure \(\mu_\gamma\) is an invariant distribution which is not represented by a smooth function.
\end{lemma}

\begin{proof}
Let \(T>0\) be a period of \(\gamma\), and parametrize \(\gamma\) by
\[
p:[0,T]\to M,\qquad p(t)=g_t(p(0)),\qquad p(0)=p(T).
\]
Define
\[
\mu_\gamma(\varphi)=\frac1T\int_0^T \varphi(p(t))\,dt,
\qquad \varphi\in C_c^\infty(M).
\]
Then \(\mu_\gamma\) is a probability measure supported on \(\gamma\). For any \(\varphi\),
\[
\langle L_X\mu_\gamma,\varphi\rangle
=
-\frac1T\int_0^T X\varphi(p(t))\,dt
=
-\frac1T\int_0^T \frac{d}{dt}\varphi(p(t))\,dt
=
-\frac{\varphi(p(T))-\varphi(p(0))}{T}
=0.
\]
Thus \(L_X\mu_\gamma=0\). Equivalently, \(\mu_\gamma\) is invariant under the flow.

Suppose now that \(\mu_\gamma\) were represented by a smooth function \(h\). Then for some positive smooth density \(\nu\),
\[
\mu_\gamma=h\nu.
\]
Since \(\gamma\neq M\) and \(\gamma\) is a compact embedded submanifold of dimension \(0\) or \(1\), it has empty interior in \(M\). Hence \(M\setminus\gamma\) is a nonempty open dense subset, and \(\mu_\gamma(M\setminus\gamma)=0\). Therefore
\[
\int_{M\setminus\gamma} h\,d\nu=0.
\]
Since \(h\) is continuous and \(\nu>0\), this forces \(h=0\) on \(M\setminus\gamma\). By density of \(M\setminus\gamma\), continuity gives \(h=0\) on all of \(M\). Thus \(\mu_\gamma=0\), contradiction. Therefore \(\mu_\gamma\) is not represented by a smooth function.
\end{proof}

\section{Positive entropy forces a periodic orbit}

We use the following standard consequence of Katok's horseshoe theorem.

\begin{theorem}[Katok]
Let \(N\) be a compact smooth manifold, and let \(F:N\to N\) be a \(C^{1+\alpha}\) diffeomorphism, \(\alpha>0\). If the topological entropy of \(F\) is positive, then \(F\) has a hyperbolic horseshoe; in particular, \(F\) has a periodic point.
\end{theorem}

For flows, the topological entropy satisfies the standard normalization
\[
h_{\mathrm{top}}(g)=h_{\mathrm{top}}(g_1).
\]
Equivalently, for every \(t\neq 0\),
\[
h_{\mathrm{top}}(g_t)=|t|\,h_{\mathrm{top}}(g).
\]
Thus positive entropy of the flow is equivalent to positive entropy of the time-one map \(g_1\).

\begin{lemma}
If \(h_{\mathrm{top}}(g)>0\), then the flow \(g_t\) has a periodic orbit.
\end{lemma}

\begin{proof}
Let \(F=g_1\). Since \(h_{\mathrm{top}}(g)>0\), we have \(h_{\mathrm{top}}(F)>0\). The map \(F\) is a \(C^\infty\) diffeomorphism of the compact manifold \(M\). By Katok's theorem, \(F\) has a periodic point: there exist \(x\in M\) and \(n\ge 1\) such that
\[
F^n(x)=g_n(x)=x.
\]
Hence \(x\) lies on a periodic orbit of the flow \(g_t\). If \(g_t(x)=x\) for all \(t\), this is the degenerate periodic orbit consisting of a fixed point.
\end{proof}

\section{Proof of the theorem}

Assume, for contradiction, that
\[
h_{\mathrm{top}}(g)>0.
\]
By the previous lemma, there exists a periodic orbit \(\gamma\) of \(X\).

Since \(h_{\mathrm{top}}(g)>0\), \(M\) cannot be a point or a circle: every continuous flow on a compact \(0\)- or \(1\)-dimensional manifold has zero topological entropy. Hence \(\dim M\ge 2\). Therefore the periodic orbit \(\gamma\), being of dimension \(0\) or \(1\), satisfies
\[
\gamma\neq M.
\]
By the orbital-measure lemma, the normalized measure \(\mu_\gamma\) on \(\gamma\) is an invariant distribution:
\[
L_X\mu_\gamma=0.
\]
Taking \(f=0\), which is a smooth function, the hypothesis implies that \(\mu_\gamma\) must be represented by a smooth function. But the orbital-measure lemma says that this is impossible because \(\gamma\neq M\). This is a contradiction.

Therefore
\[
h_{\mathrm{top}}(g)=0.
\]
Thus the flow generated by \(X\) has zero topological entropy.

\end{document}